\documentclass[11pt,a4paper]{article}
\usepackage[T1]{fontenc}
\usepackage[utf8]{inputenc}
\usepackage{lmodern}
\usepackage[a4paper,margin=1in]{geometry}
\usepackage{microtype}
\usepackage{parskip}
\usepackage{enumitem}
\usepackage[hidelinks]{hyperref}

\setlist[itemize]{leftmargin=1.5em,itemsep=0.3em,topsep=0.3em}
\setlength{\emergencystretch}{2em}
\pagestyle{empty}

\begin{document}

\begin{center}
{\Large\bfseries Cover Letter: Response to Reviewers}\\[0.6em]
{\large Submission 241}\\[0.3em]
{\large Benchmark for Crack Re-Identification in Ticket Management Systems}
\end{center}

\vspace{0.3em}

We thank both reviewers for their positive assessments and careful reading.
Neither review requested specific changes. Their main concerns, namely the small
dataset (11 test walls and 19 matchable identities), single-visit acquisition,
single-annotator labels, and partial registration coverage, are the central
limitations of the study and are stated in Section~VIII, the Abstract, and the
Conclusion. New walls, visits, and annotations could not be collected within the
camera-ready period, so no experiments were re-run and all reported results are
unchanged. The revision instead makes the preliminary scope of the findings
clearer, as Reviewer~2 asked:

\begin{itemize}
  \item \textbf{Claims calibrated to the evidence.} Validated F1 now ``stays
  close to'' rather than ``slightly exceeds'' the constant baseline; ``outperform
  chance'' is restricted to Rank-1; and statements about wall context and
  geometric alignment are phrased as suggestive (``suggest'', ``may'') rather
  than conclusive.
  \item \textbf{Corrections.} The registration gate is now correctly described
  as retaining 36.8\% of negatives (not discarding them), and the abstract's
  ``twice'' comparison is replaced by the two DIR values (0.768 and 0.429), with
  a note that they are computed over different candidate pools.
  \item \textbf{Clarifications from existing data and code.} We added the
  breakdown of the unclaimed annotation points, the synthetic-warp parameters, a
  definition of HYBRID-SKEL, and a citation for the inpainting method~[39]; we
  also clarified the purpose of the sharpness threshold and the construction of
  the reference-level baselines.
  \item \textbf{Consistency.} Terminology was harmonized (registered pairs, valid
  pairs, coverage, within-visit versus cross-visit), and citations and
  cross-references were verified.
\end{itemize}

\section*{Reviewer 1}

We thank the reviewer for the positive assessment of the problem, the dataset,
and the reporting of negative results. We agree that a larger dataset and a real
cross-visit evaluation would strengthen the contribution; both are the main
directions for future work (Section~VIII). The paper is positioned as an initial
benchmark and a negative result, and the current results are not presented as
cross-visit evidence.

\section*{Reviewer 2}

We thank the reviewer for the assessment of the open-set protocol and of the
context and geometry analyses. We agree that the small sample, single-annotator
labels, limited registration coverage, and absence of cross-visit evaluation bound
the generalizability of the conclusions. In line with the reviewer's condition
that the preliminary nature of the findings be clearly communicated, we toned down
the claims listed above and kept the point-estimate and single-visit caveats in
the Abstract, Section~VIII, and the Conclusion.

\vspace{0.5em}

We thank the reviewers again for their time and feedback.

\medskip

Sincerely,\\
The authors

\end{document}